\documentclass[11pt]{article}
\usepackage[margin=1in]{geometry}
\usepackage{amsmath, amssymb}
\usepackage{xcolor}
\usepackage{enumitem}
\usepackage{listings}
\usepackage{hyperref}

\lstset{
  basicstyle=\ttfamily\small,
  columns=fullflexible,
  keepspaces=true,
  breaklines=true,
  frame=single,
  xleftmargin=1em,
  framexleftmargin=1em,
}

\newcommand{\code}[1]{\texttt{#1}}
\newcommand{\todo}[1]{\textcolor{red}{[#1]}}

\title{Feature-Pair Refactor of the Polygon Force Kernels\\
  \large (biarc target; normal and single-arc rounded as stepping stones)}
\author{}
\date{}

\begin{document}
\maketitle

\section{Scope and terminology}

The physical target model is the \textbf{biarc} polygon:\ each vertex
has two tangent-continuous circular arcs meeting at a junction $j$ on
the inward bisector (\S 4 of \emph{packing\_treatise-7}). Polygon $s$
with $n_s$ vertices therefore has
\begin{itemize}[leftmargin=*]
  \item $n_s$ straight edges (one per consecutive vertex pair), and
  \item $2 n_s$ arcs (two per vertex: ``arc 1'' from $a_i^-$ to the
        junction $j_i$, ``arc 2'' from $j_i$ to $a_i^+$).
\end{itemize}
So $3 n_s$ features total per polygon.

The codebase currently has a \emph{single-arc rounded} kernel
(\code{updateForceEnergyRoundedKernel}) that approximates biarcs by a
single arc per vertex --- $2 n_s$ features ($n_s$ edges + $n_s$ arcs).
That kernel is FD-validated and useful, but it is a stepping stone.
The \emph{normal} model uses only edges ($n_s$ features per polygon)
and is structurally closest to what we want to refactor first.

We will build the refactor parametric in the feature set per vertex,
so the same Phase A/B/C scaffolding services all three models. The
biarc model is the destination; we walk to it via normal $\to$
single-arc rounded $\to$ biarc.

\section{Problem statement}

The current single-arc rounded force/energy kernel
(\code{updateForceEnergyRoundedKernel}) uses per-vertex parallelism, but
each thread performs $O(n)$ sequential work inside the kernel: it loads
polygon $B$'s full geometry, computes $B$'s per-vertex arc geometry,
walks all of $B$'s edges/arcs to find crossings with the thread's local
features, then walks crossings in order. With $n=32$, this serializes a
32$\times$ factor inside every thread before any of the GPU
parallelism is brought to bear. The biarc kernel would inherit and
amplify this (3$\times$ more features per vertex) if we built it the
same way.

Worse, because each thread does its own ``does my feature go inside
$B$?'' determination via $B$'s local geometry, an interior vertex deep
in $B$ (far from any edge crossing) only learns it is inside if the ball
list reaches a $B$-vertex from that interior point. That forces the
search radius up to $\sim$polygon\_radius, which is independent of the
vertex-pair interaction range. The result is a slow kernel and a
fragile, polygon-size-dependent search radius.

The same Green's-theorem accounting can be done much more cleanly --- and
the normal model already does most of it --- by reformulating around
\emph{feature pairs} and a global sorted intersection list. This document
specifies the refactor.

\section{Architecture overview}

Replace each force/energy call (normal, rounded, areaSquared) with three
phases of small per-thread kernels. Every kernel is per-pair or
per-feature, with $O(1)$ work per thread aside from inner loops bounded
by physics (number of crossings on a single feature, typically $\le 4$).

\begin{itemize}[leftmargin=*]
  \item \textbf{Phase A --- Build.} For each vertex-vertex candidate
        pair $(i_a, i_b)$ in the existing per-vertex ball list, emit the
        feature-pair candidates the chosen model requires
        (Section~\ref{sec:phaseA}).
  \item \textbf{Phase B --- Probe.} One thread per feature-pair
        candidate. Compute pairwise intersection. If there is a crossing,
        emit the crossing record. Add the per-crossing
        (``external'') atomic-add force contributions to the four
        endpoint vertices and the partner-side coupling
        (Section~\ref{sec:phaseB}).
  \item \textbf{Phase C --- Assemble.} Sort the global crossing list,
        then one thread per $A$-feature walks its sorted crossings to
        determine the inside/outside chord segments and contribute the
        per-segment (``internal'') Green's integral terms to energy and
        to the chord endpoints (Section~\ref{sec:phaseC}).
\end{itemize}

Two consequences worth highlighting up front:

\paragraph{Search radius returns to $2L$.} Because the inside/outside
state of each $A$-feature is recovered from the sorted crossing list
in Phase C --- not from each thread independently probing $B$'s geometry
--- only the crossings themselves need to be discovered. A crossing of
two convex edges of lengths at most $L$ puts the four endpoint vertices
within $2L$ of each other, so the per-vertex ball reach can be
$\sim 2.5L$ and that suffices. The polygon-size dependence in the search
radius disappears.

\paragraph{No \code{processedShapes} cap.} The current kernel dedups
neighbor shapes per thread and caps unique shapes at 32; the cap was a
performance band-aid that also subtly biased results when more shapes
existed than fit. After the refactor we enumerate \emph{feature pairs},
not neighbor shapes, and there is no fan-out cap.

\section{Feature representation}\label{sec:features}

A \emph{feature} is one segment of a polygon's boundary. We support
three feature types and bind each to a vertex index $i$:

\begin{itemize}[leftmargin=*]
  \item \code{EDGE} at $i$: the straight segment between two consecutive
        tangent points. For the normal model, this is just the segment
        from vertex $i$ to vertex $i+1$. Used by all three models.
  \item \code{ARC1} at $i$: the first arc at vertex $i$, from $a_i^-$
        to the junction $j_i$. Used by biarc only. For the single-arc
        rounded model we collapse \code{ARC1} and \code{ARC2} into one
        feature labeled \code{ARC1} (the single arc) and never emit
        \code{ARC2}.
  \item \code{ARC2} at $i$: the second arc, from $j_i$ to $a_i^+$. Used
        by biarc only.
\end{itemize}

We encode a feature ID in a single \code{uint32\_t} with the type in
the low two bits:
\begin{equation*}
  \text{featureId}(s, i, t)
  \;=\; 4 \cdot (\text{startIndices}[s] + i) + t,
  \qquad t \in \{\text{EDGE}=0,\; \text{ARC1}=1,\; \text{ARC2}=2\}.
\end{equation*}
The global vertex index is $\lfloor \text{featureId} / 4 \rfloor$,
the type is \code{featureId \& 3}, and \code{shapeId[\,]} gives the
owning polygon.

Feature counts per polygon by model:
\begin{center}
\begin{tabular}{lrr}
  Model & features per vertex & features per polygon \\
  \hline
  normal             & 1 (\code{EDGE})                          & $n$  \\
  single-arc rounded & 2 (\code{EDGE}, \code{ARC1})             & $2n$ \\
  biarc              & 3 (\code{EDGE}, \code{ARC1}, \code{ARC2})& $3n$ \\
\end{tabular}
\end{center}

\section{Phase A: candidate generation}\label{sec:phaseA}

\textbf{Input.} The per-vertex ball list
(\code{ballNeighbors[\,]}, \code{numBallNeighbors[\,]}) already
populated by either \code{updateNeighborCells()} or
\code{updateNeighborBall()}. This phase reuses Decision 1's choice ---
the vertex-vertex list is the source of truth for spatial proximity.

\textbf{Output.} A flat candidate-pair array
\code{featurePairs[\,]} where each entry is
\begin{lstlisting}
struct FeaturePair { uint32_t fA; uint32_t fB; };
\end{lstlisting}
plus a count.

\textbf{Logic.} One thread per entry of the per-vertex ball list. For
the pair $(i_a, i_b)$ that thread is handling, it emits, depending on
the active model:

For each vertex-vertex candidate $(i_a, i_b)$ we emit the Cartesian
product of $i_a$'s feature set with $i_b$'s feature set:

\begin{center}
\begin{tabular}{lrl}
  Model              & pairs per $(i_a, i_b)$ & types emitted \\
  \hline
  normal             & $1 \times 1 = 1$  & edge$\times$edge \\
  single-arc rounded & $2 \times 2 = 4$  & all combinations of \{edge, arc1\} \\
  biarc              & $3 \times 3 = 9$  & all combinations of \{edge, arc1, arc2\} \\
\end{tabular}
\end{center}

Per the user's direction (``assume the worst''), we emit \emph{all}
pairs unconditionally, including the (arc, arc) combinations where a
crossing is rare. We accept the extra memory in exchange for not
having to reason about edge-case filtering.

The candidate count grows by a factor of 9 (biarc) over the
vertex-vertex list. With $N \cdot n = 2048$ vertices and a typical 575
candidates per vertex (Section~\ref{sec:perf-estimates}), this is
roughly $9 \cdot 2048 \cdot 575 / 2 \approx 5.3 \times 10^6$
pairs at peak. Stored as a pair of \code{uint32\_t} per entry, that's
$\sim$42~MB at peak in the worst case (pre-relax, all-to-all);
post-relax it drops by $\sim$200$\times$.

The emission is parallel and produces output to consecutive locations.
We allocate the output array with a worst-case bound
($4 \cdot \text{ballMaxNeighbors} \cdot \text{size}$) and use a single
\code{atomicAdd} on a shared counter to claim slots, or use CUB's
\code{DeviceScan} to compact in a second pass --- whichever is
simpler. Resizing on overflow follows the existing pattern.

\section{Phase B: pairwise probe}\label{sec:phaseB}

\textbf{Input.} \code{featurePairs[\,]} from Phase A.

\textbf{Output.}
\begin{itemize}[leftmargin=*]
  \item A flat array of crossing records (one per crossing, with
        $A$-feature ID and parameter for sorting).
  \item Direct \code{atomicAdd}-style force contributions for the
        ``external'' per-crossing terms.
\end{itemize}

\begin{lstlisting}
struct Crossing {
    uint32_t fA;          // owning A-feature
    uint32_t fB;          // partner B-feature
    float    paramA;      // parameter along fA at crossing (for sorting)
    double   Xx, Xy;      // crossing point
    double   eBx, eBy;    // tangent direction of fB at crossing (for forces)
    double   paramB;      // parameter along fB at crossing
};
\end{lstlisting}

\textbf{Kernel.} One thread per entry of \code{featurePairs[\,]}. The
thread reads the pair $(f_a, f_b)$, decodes their types (low bit),
loads the four (or six --- for arcs, plus centers and angles) coordinates
needed, and dispatches to the appropriate pairwise intersection
routine:

\begin{itemize}[leftmargin=*]
  \item \code{edgeEdgeCross} --- two line-segments. 0 or 1 crossing.
        (Already exists in the codebase as part of
        \code{updateNeighborsKernel}; factor it into
        \code{intersect.cuh}.)
  \item \code{edgeArcCross} --- segment intersect circle, gated to the
        arc parameter range and the edge range. 0, 1, or 2 crossings.
        Works identically for \code{ARC1} and \code{ARC2}; the only
        per-arc data passed in are center, radius, and angle range.
  \item \code{arcEdgeCross} --- same as above with the roles swapped.
  \item \code{arcArcCross} --- two circles. 0, 1, or 2 crossings. Also
        agnostic to whether each side is \code{ARC1} or \code{ARC2}.
\end{itemize}

The dispatch in Phase B's kernel decodes the type bits of the two
feature IDs to pick the right routine, and reads the appropriate
center/radius/range data for each arc type. Biarc adds no new
intersection math --- just two more arc variants on each side.

Each crossing is appended to the output array (atomic counter). For
each crossing, the thread also applies the external force
contributions: the chain-rule force from the implicit function theorem
($dX_c/dv$ contributions on the four endpoint vertices and the
partner-feature coupling on $B$'s side). These are local to the
crossing and do not depend on inside/outside state, so they fire
here.

\textbf{No inner loops over $n$ anywhere in this phase.}

\section{Phase C: assembly}\label{sec:phaseC}

\textbf{Step 1 --- global sort.} Sort the crossing records by
\code{(fA, paramA)} using CUB's \code{DeviceRadixSort::SortPairs}, the
same pattern \code{updateOutersections} already uses for the normal
model. The packed sort key can be \code{(fA << 32) | floatBits(paramA)}
into a \code{uint64\_t}, then a permutation applied to the full record
array.

\textbf{Step 2 --- per-feature walk.} One thread per $A$-feature reads
its contiguous slice of the sorted list, sorts the crossings by
\code{paramA} (at most $\sim$4 entries; a small insertion sort), and
walks them in order:

\begin{enumerate}[leftmargin=*]
  \item Determine the inside/outside state at the start of the feature.
        For \code{EDGE} at vertex $i$, this is ``is vertex
        $i$ inside any of the polygons whose features cross this edge?''
        Equivalently: walk the polygon-pair list of this $A$-feature,
        count crossings up to the start parameter, and take parity.
        \emph{Implementation note:} since each polygon $B$ that
        crosses $f_a$ contributes an even number of crossings on $f_a$
        when $f_a$ enters/exits $B$ completely, the parity at the start
        can be deduced from how many $B$-polygons contain $f_a$'s start
        endpoint. This is the same ``which shapes contain me'' logic
        the normal model already implements.
  \item For each consecutive pair of crossings (and the segments
        before the first and after the last), determine whether that
        sub-segment of $f_a$ is inside the relevant other polygon, and
        if so, contribute the internal Green's-theorem term for that
        sub-segment to:
        \begin{itemize}
          \item the global energy (overlap area term), and
          \item the forces on the segment's endpoint vertices
                (the chord-line-integral force pieces).
        \end{itemize}
\end{enumerate}

For \code{ARC} features the integration is over a circular arc rather
than a straight segment; the per-segment integrand is the standard
arc-area-with-respect-to-origin formula. Edge features use the
straight-line variant.

\textbf{No inner loops over $n$ anywhere in this phase} either ---
the walk is over the feature's own crossings, bounded by physics.

\section{Shared helpers}\label{sec:helpers}

Per Decision 5, factor out the routines that are used across normal,
rounded, and areaSquared:

\begin{itemize}[leftmargin=*]
  \item \code{src/features.cuh} --- Feature ID encode/decode,
        per-feature geometry accessors (edge endpoints, arc center +
        radius + angle range, all derivable from \code{positions[\,]}
        and \code{shapeId/startIndices/next/prev}).
  \item \code{src/intersect.cuh} --- the four pairwise intersection
        routines (\code{edgeEdgeCross}, \code{edgeArcCross},
        \code{arcEdgeCross}, \code{arcArcCross}). Each returns its
        crossings as a small fixed-size buffer.
  \item \code{src/integrate.cuh} --- per-segment Green's contribution
        for an edge chord and an arc chord, with respect to a per-pair
        anchor. Force gradients of these terms.
  \item \code{src/parityWalk.cuh} --- the ``walk sorted crossings and
        figure out which shapes contain me'' helper. Used identically
        by the normal model's existing flow and by Phase~C above.
  \item \code{src/forces/normal.cu} --- normal model entry point:
        Phase A emits edge$\times$edge only, then B and C.
  \item \code{src/forces/rounded.cu} --- single-arc rounded entry
        point: Phase A emits the 4 \{edge, arc1\} combinations.
  \item \code{src/forces/biarc.cu} --- biarc entry point: Phase A
        emits the 9 \{edge, arc1, arc2\} combinations.
  \item \code{src/forces/areaSquared.cu} --- wraps biarc (or rounded)
        with the existing 2-pass area-scaling.
\end{itemize}

\section{Implementation order}\label{sec:order}

Following Decision 5: start with the normal model. Then walk up the
feature-set ladder.

\begin{enumerate}[leftmargin=*]
  \item \textbf{Phase 0 --- shared scaffolding (no behavior change).}
        Create the headers above with the encode/decode and helper
        signatures; move the existing line-line intersection logic out
        of \code{updateNeighborsKernel} into \code{intersect.cuh}. Verify
        normal model still passes existing FD tests.
  \item \textbf{Phase 1 --- normal model on new architecture.}
        Replace \code{updateNeighbors / updateOutersections /
        updateForceEnergyExterior / updateForceEnergyInterior /
        updateForceEnergyVertexDisk} with the new A/B/C flow. The
        existing logic is conceptually already this shape, so this
        step is mostly restructuring + factoring shared helpers.
        Phase A emits only edge$\times$edge.
        Validate with the existing normal-model FD tests: force =
        $-\nabla E$ to the existing 0.1\% tolerance.
  \item \textbf{Phase 2 --- single-arc rounded model on new
        architecture.} Add the line-arc and arc-arc routines to
        \code{intersect.cuh} and the arc-segment Green's terms to
        \code{integrate.cuh}. Phase A emits the 4 combinations of
        \{edge, arc1\}. Build \code{forces/rounded.cu}. Validate with
        the existing rounded FD tests: force = $-\nabla E$ to the
        existing $1.6 \times 10^{-5}$ tolerance.
  \item \textbf{Phase 3 --- biarc model.} Add \code{ARC2} as a feature
        type and the biarc-junction geometry (radii $d_1, d_2$,
        centers $z_1, z_2$, angle ranges from \S 4 of
        \emph{packing\_treatise-7}). Phase A emits the 9 combinations
        of \{edge, arc1, arc2\}. Phase B's intersection dispatch
        already covers (arc, arc) and (arc, edge); no new math, only
        the geometry-construction helper needs an \code{arc2} variant.
        Phase C's arc-segment Green's term applies unchanged.
        FD-validate against the same target tolerance as rounded.
  \item \textbf{Phase 4 --- areaSquared.} Two-pass mode on top of
        biarc (or rounded, depending on what's wanted). Validate
        with existing areaSquared tests.
  \item \textbf{Phase 5 --- delete the old code paths.} See
        Section~\ref{sec:delete}.
\end{enumerate}

Each phase is independently validated. The old per-vertex
\code{updateForceEnergyRoundedKernel} stays compileable in the source
tree until Phase 4 so we have a reference to FD against if Phase 2
introduces a regression.

\section{What we delete after Phase 4}\label{sec:delete}

Per Decision 6 (and confirmed: anything not compatible):

\begin{itemize}[leftmargin=*]
  \item \code{updateForceEnergyRoundedKernel} (the 1500-line monolith).
  \item \code{updateForceEnergyAreaSquaredCUDA} (replaced by the new
        2-pass wrapper).
  \item The per-thread \code{roundedScratch} buffer of $38n$ doubles
        per vertex.
  \item The \code{processedShapes[32]} dedup logic and the
        \code{maxProc} cap.
  \item The polygon-aware \code{searchFactor} default in
        \code{setStartIndices} (revert to fixed 2.5).
  \item The polygon-aware \code{trueCutoff} in
        \code{updateNeighborBall} (revert to $2 \cdot
        \text{maxEdgeLength}$).
  \item The \code{updateForceEnergyExteriorKernel} and
        \code{updateForceEnergyInteriorKernel} (their work is absorbed
        into Phase B and Phase C respectively).
  \item \code{updateForceEnergyVertexDiskKernel} and its ball variant
        --- the vertex-disk cap is just a special case of
        edge-edge / arc-edge crossings in the new flow.
  \item The compatibility scaffolding around two neighborTypes can
        stay; both \code{updateNeighborCells} and
        \code{updateNeighborBall} keep populating the per-vertex
        candidate buffer that Phase A reads.
\end{itemize}

\section{Memory and performance estimates}\label{sec:perf-estimates}

For the $N{=}64$, $n{=}32$, $\phi{=}1.1$ working case:

\begin{tabular}{lr}
  Total vertices ($N n$): & 2048 \\
  Candidates per vertex (post-relax, $\text{ball}=2.5L$): & $\sim$12 \\
  Candidates per vertex (pre-relax, all-to-all): & 2016 \\
  Feature-pairs (normal, peak): & $2048 \cdot 2016 / 2 \approx 2 \times 10^6$ \\
  Feature-pairs (rounded, peak): & $4 \cdot 2 \times 10^6 = 8 \times 10^6$ \\
  Feature-pairs (biarc, peak): & $9 \cdot 2 \times 10^6 \approx 1.8 \times 10^7$ \\
  Feature-pair memory at biarc peak: & $\sim$144~MB (one shot, pre-relax) \\
  Feature-pair memory post-relax: & $\sim$0.7~MB \\
  Crossings emitted (rough estimate): & $\sim 10^5$ \\
  Crossings memory: & $\sim$5~MB \\
\end{tabular}

After softBody relaxation (when polygons are at target shapes and
ball list is sparse), the feature-pair count drops by $\sim 200\times$
to $\sim$40K, well within memory budgets and dominated by Phase B's
$O(1)$-per-thread work.

\section{Open items for later}

\begin{itemize}[leftmargin=*]
  \item Containment cases (one polygon fully inside another with no
        edge crossings) emit zero crossings from Phase B. Phase C
        handles this correctly \emph{if} we seed it with a polygon-pair
        ``A's start endpoint inside B?'' flag. Decide whether that
        flag comes from a separate cheap kernel or from a sentinel
        crossing emitted by Phase A. Punt to Phase 1 implementation.
  \item Whether to express Phase A's emission via a template parameter
        on \code{modelType} (compile-time) or an enum branch (runtime).
        Compile-time is cleaner; we can do either.
\end{itemize}

\end{document}
